% !TEX root = ../main.tex
\section{Three Fundamental Estimators and Similarity}
\label{sec:algorithms-and-errors}

\subsection{SGD, MSGD, and TSGD}
\label{subsec:algorithms-and-assumptions}

\begin{assumption}[Smooth, lower-bounded objective]
\label{ass:smooth-objective}
The function $f:\mathbb R^d\to\mathbb R$ has an $L$-Lipschitz gradient,
with $L>0$, and $f_{\inf}:=\inf_x f(x)>-\infty$.
\end{assumption}
\begin{assumption}[Unbiased stochastic gradients]
\label{ass:unbiased-oracle}
The oracle satisfies $\mathbb E_{\xi}[h(x,\xi)]=\nabla f(x)$ for every $x$.
Each new sample is independent of the preceding history.
\end{assumption}

All three methods use positive deterministic stepsizes $\eta_k$ and
the same iterate update, starting from a deterministic $x_0$:
\begin{equation}
\label{eq:tsgd-iterate}
x_{k+1}=x_k-\eta_k g_k.
\end{equation}
They differ in how they construct $g_{k+1}$, their estimate of
$\nabla f(x_{k+1})$:
\begin{align}
\mathrm{SGD}:\quad g_{k+1}
  &=h(x_{k+1},\xi_{k+1}), \label{eq:sgd-estimate}\\
\mathrm{MSGD}:\quad g_{k+1}
  &=(1-\beta_k)g_k+\beta_k h(x_{k+1},\xi_{k+1}),
  \quad 0<\beta_k\leq1, \label{eq:msgd-estimate}\\
\mathrm{TSGD}:\quad g_{k+1}
  &=g_k+D_{k+1}. \label{eq:tsgd-estimate}
\end{align}
In the standard pointwise-oracle construction,
\begin{equation}
\label{eq:shared-sample-increment}
D_{k+1}=h(x_{k+1},\xi_{k+1})-h(x_k,\xi_{k+1}).
\end{equation}
Both evaluations use the \emph{same fresh sample}. SGD draws a fresh
gradient estimate, MSGD averages it with the previous estimate, and TSGD
adds an estimated gradient change from $x_k$ to $x_{k+1}$ to $g_k$,
transporting it to the new iterate. Momentum averaging can reduce
noise, but gradients from earlier iterates can make the estimate lag
behind the current gradient, especially when the gradient changes rapidly.
Transport accounts for this change but inherits the estimation error
in $g_k$. Unless specified otherwise, SGD and MSGD initialize with
$g_0=h(x_0,\xi_0)$. TSGD permits any $g_0$ with finite mean-square
error, including a biased estimate.

\subsection{Understanding Similarity}
\label{subsec:understanding-similarity}

Write $r(x,\xi)=h(x,\xi)-\nabla f(x)$ for the gradient noise.
\begin{assumption}[Bounded variance (BV)]
\label{ass:bounded-variance}
For some $\sigma\geq0$,
\begin{equation}
\label{eq:bounded-variance}
\mathbb E_{\xi}[\|r(x,\xi)\|^2]\leq\sigma^2
\qquad\text{for all }x.
\end{equation}
\end{assumption}
Transport instead uses a condition on gradient differences.
\begin{assumption}[Centered $\delta$-similarity]
\label{ass:oracle-similarity}
For some $\delta\geq0$,
\begin{equation}
\label{eq:oracle-similarity}
\mathbb E_{\xi}[\|r(x,\xi)-r(y,\xi)\|^2]
\leq\delta^2\|x-y\|^2
\qquad\text{for all }x,y.
\end{equation}
\end{assumption}
Here $\sigma$ bounds the magnitude of individual gradient errors, whereas
$\delta$ measures their variation under shared sampling. A small $\delta$
means that stochastic gradient differences closely match deterministic
gradient differences; it does not require the gradients themselves, or
their changes, to be small. Similarity is a condition on the stochastic
oracle and applies to both finite-sum and population objectives.

For an illustration, let $f=n^{-1}\sum_i f_i$ with
$f_i(x)=\tfrac12x^\top A_i x+b_i^\top x+c_i$ and symmetric $A_i$.
Uniformly sampling $i$, and writing $\bar A=n^{-1}\sum_i A_i$, gives
\begin{equation}
\label{eq:quadratic-difference-noise}
[\nabla f_i(x)-\nabla f_i(y)]-[\nabla f(x)-\nabla f(y)]
=(A_i-\bar A)(x-y).
\end{equation}
The linear terms $b_i^\top x$ cancel from the gradient differences.
Components can therefore have very different
gradients but nearly identical gradient changes. Identical Hessians give
exact transport even with arbitrarily large additive gradient noise.
Hessian heterogeneity controls $\delta$, whereas $L$ concerns the average
Hessian; neither must be smaller than the other
(Appendix~\ref{app:noise-assumptions}).

Scalar least squares separates the two quantities. Let
$f(x)=\tfrac12\mathbb E[(ux-b)^2]$, with feature $u$ satisfying
$0<\mathbb E[u^2]$ and $\mathbb E[u^4]<\infty$, and label
$b=ux^\star+\varepsilon$. The label noise $\varepsilon$ is independent
of $u$, with $\mathbb E[\varepsilon]=0$ and
$\mathbb E[\varepsilon^2]=s^2<\infty$. The stochastic gradient is
$h(x;u,\varepsilon)=u^2(x-x^\star)-u\varepsilon$, so
\begin{equation}
\label{eq:least-squares-similarity}
\begin{aligned}
h(x;u,\varepsilon)-h(y;u,\varepsilon)&=u^2(x-y),\\
\mathbb E[|h(x;u,\varepsilon)-\nabla f(x)|^2]
&=\delta^2(x-x^\star)^2+Ls^2,
\end{aligned}
\end{equation}
where $\delta^2=\operatorname{Var}(u^2)$ and $L=\mathbb E[u^2]$.
The two terms play different roles. Label noise contributes to the noise
magnitude but cancels from differences, while feature variability sets
$\delta$. Increasing $s$ therefore makes every fresh gradient noisier
without affecting transport. When $\delta>0$, the pointwise noise
variance is unbounded in $x$, so similarity holds but global BV fails.
Label noise still affects TSGD through its initial estimate, but not
through later transport increments. Sections~\ref{subsec:additive-noise-comparison}
and~\ref{subsec:multiplicative-noise-comparison} illustrate the two
extremes of this model: a constant feature ($\delta=0$) with label noise
favors transport with accurate initialization, whereas noiseless labels
($s=0$) with varying features make fresh gradients exact at the minimizer.

\subsection{Two Fundamental Sources of TSGD Error}
\label{subsec:martingale-interpretation}

Define $e_k=g_k-\nabla f(x_k)$, $E_k=\mathbb E[\|e_k\|^2]$, and
$d_{k+1}=D_{k+1}-[\nabla f(x_{k+1})-\nabla f(x_k)]$.
The transport update gives
\begin{equation}
\label{eq:estimation-error-recursion}
\boxed{e_{k+1}=e_k+d_{k+1}.}
\end{equation}
For conditionally unbiased increments, the new error has zero conditional
mean and is orthogonal in expectation to previous errors. Consequently,
\begin{equation}
\label{eq:transport-error-identity}
E_k=E_0+\sum_{t=0}^{k-1}\mathbb E[\|d_{t+1}\|^2].
\end{equation}
Initialization error persists, while transport errors accumulate.
Similarity bounds each contribution by the squared displacement, yielding
\begin{equation}
\label{eq:transport-error-growth}
E_k\leq E_0+\delta^2\sum_{t=0}^{k-1}
\eta_t^2\mathbb E[\|g_t\|^2].
\end{equation}
This explains both the opportunity and the difficulty: accurate changes
preserve useful information, but even small errors can accumulate.
Section~\ref{sec:tsgd-convergence} states the conditional assumptions
directly on $D_{k+1}$, allowing transport beyond pointwise gradient oracles.
